% ============================================================
% INFORME FINAL -- CEDENAR
% Universidad Cooperativa de Colombia, Campus Pasto
% Análisis de Software -- Semestre III
% Formato: APA 7
%
% INSTRUCCIONES PARA OVERLEAF:
% Subir las siguientes imágenes antes de compilar:
%   - diagrama_clases.png      (wiki, sección 3.1)
%   - diagrama_paquetes.png    (wiki, sección 3.2)
%   - diagrama_componentes.png (wiki, sección 3.3)
%   - diagrama_despliegue.png  (wiki, sección 3.4)
%   - diagrama_casos_uso.png   (wiki, sección 4.1)
%   - diagrama_actividades.png (wiki, sección 4.2)
%   - diagrama_secuencia.png   (wiki, sección 4.3)
%   - maquina_estados.png      (wiki, sección 4.4)
%   - vista_mobile.png         (archivo local: diagramas/vista mobile.png)
%   - archimate.png            (diagrama ArchiMate, sección 2)
% Compilar con pdfLaTeX
% ============================================================

\documentclass[12pt, a4paper]{article}

% --- Codificación y lenguaje ---
\usepackage[utf8]{inputenc}
\usepackage[T1]{fontenc}
\usepackage[spanish, es-tabla]{babel}

% --- Fuente Times New Roman (APA 7) ---
\usepackage{mathptmx}

% --- Márgenes APA 7: 2.54 cm todos los lados ---
\usepackage[top=2.54cm, bottom=2.54cm, left=2.54cm, right=2.54cm]{geometry}

% --- Interlineado 2.0 (APA 7) ---
\usepackage{setspace}
\doublespacing

% --- Imágenes ---
\usepackage{graphicx}
\usepackage{float}

% --- Tablas ---
\usepackage{booktabs}
\usepackage{longtable}
\usepackage{array}
\usepackage{tabularx}

% --- Código fuente sin color ---
\usepackage{listings}

\lstset{
  basicstyle=\ttfamily\small,
  breaklines=true,
  breakatwhitespace=true,
  frame=single,
  xleftmargin=12pt,
  xrightmargin=6pt,
  aboveskip=6pt,
  belowskip=6pt,
  numbers=left,
  numberstyle=\tiny,
  numbersep=8pt,
  tabsize=4,
  showstringspaces=false,
  captionpos=b,
  keepspaces=true
}

% --- Párrafos APA 7 ---
\setlength{\parindent}{1.27cm}
\setlength{\parskip}{0pt}

% --- Encabezado APA 7 ---
\usepackage{fancyhdr}
\pagestyle{fancy}
\fancyhf{}
\fancyhead[L]{\small SISTEMA DE AUTOMATIZACIÓN DE LECTURAS -- CEDENAR}
\fancyhead[R]{\small \thepage}
\renewcommand{\headrulewidth}{0pt}

% --- Secciones APA 7: centradas, negrita, nivel 1 ---
\usepackage{titlesec}
\titleformat{\section}{\normalfont\bfseries\centering}{}{0em}{}
\titleformat{\subsection}{\normalfont\bfseries}{}{0em}{}
\titleformat{\subsubsection}{\normalfont\bfseries\itshape}{}{0em}{}
\titleformat{\paragraph}[runin]{\normalfont\bfseries}{}{0em}{}[.]

% --- Links sin color ---
\usepackage{hyperref}
\hypersetup{
  colorlinks=false,
  pdfborder={0 0 0},
  pdftitle={Sistema de Automatización de Lecturas CEDENAR},
  pdfauthor={Urbano, Salazar, Diaz, Parra, Lopez}
}

% --- Captions APA 7 ---
\usepackage{caption}
\captionsetup{
  labelfont=bf,
  textfont=it,
  labelsep=newline,
  justification=raggedright,
  singlelinecheck=false,
  aboveskip=4pt
}

% ============================================================
\begin{document}

% ============================================================
% PORTADA APA 7
% ============================================================
\begin{titlepage}
  \centering
  \vspace*{3cm}

  {\large\bfseries
    Sistema de Automatización de Lecturas de Medidores de Energía\\
    para las Centrales Eléctricas de Nariño (CEDENAR)
  \par}

  \vspace{1.5cm}

  {\normalsize
    Jenifer Daniela Urbano Córdoba, Drako David Salazar Torres,\\
    Nicolas Alejandro Diaz Acosta, Cristhian Santiago Parra Erazo\\
    y Juan Camilo Lopez Diaz
  \par}

  \vspace{1cm}

  {\normalsize
    Programa de Ingeniería de Software, Semestre III\\
    Universidad Cooperativa de Colombia, Campus Pasto
  \par}

  \vspace{1cm}

  {\normalsize
    Análisis de Software\\
    Oscar Andrés Osorio
  \par}

  \vspace{1cm}

  {\normalsize 25 de mayo de 2026\par}

\end{titlepage}

% ============================================================
% RESUMEN
% ============================================================
\newpage
\begin{center}
  \textbf{Resumen}
\end{center}

\noindent Este informe documenta el análisis y diseño de un sistema de software para automatizar la toma de lecturas de medidores de energía eléctrica en las Centrales Eléctricas de Nariño (CEDENAR). El proceso actual presenta tres problemas concretos: dependencia de conectividad continua durante las jornadas de campo, ausencia de canal formal para que los clientes reporten fallas o aclaren cobros, y uso de papel e impresora manual para los comprobantes. La solución propuesta es una aplicación móvil con arquitectura \textit{offline-first} para el operario de campo y un portal web para el cliente. El análisis cubre levantamiento de requerimientos funcionales y no funcionales, seis historias de usuario con criterios de aceptación en Gherkin, modelado de arquitectura empresarial con ArchiMate, seis diagramas UML estáticos y dinámicos, aplicación de la metodología Design Thinking con benchmarking de soluciones existentes, un prototipo de alta fidelidad en Figma y fragmentos de código en Kotlin que ilustran la lógica de negocio central.

\vspace{6pt}
\noindent\textbf{Palabras clave:} automatización de lecturas, arquitectura \textit{offline-first}, análisis de software, UML, Design Thinking.

% ============================================================
\newpage
\tableofcontents
\newpage

% ============================================================
\section{Introducción y Levantamiento de Requerimientos}
% ============================================================

\subsection{Contexto del Problema}

\noindent CEDENAR opera en el departamento de Nariño con recaudadores instaladores que recorren zonas urbanas y rurales tomando lecturas de medidores de energía eléctrica. El proceso depende de una terminal portátil que descarga los datos del sector antes de cada jornada y requiere conexión activa al servidor para operar. Cuando la señal cae o la pantalla del dispositivo se refresca de manera accidental, el equipo puede congelarse hasta 15 minutos y el operario pierde el avance registrado hasta ese momento.

Del lado del cliente residencial, no existe canal formal para aclarar cobros ni reportar interrupciones del servicio. La única alternativa disponible es acudir directamente a la residencia del recaudador instalador del sector o enviar un mensaje por WhatsApp personal al técnico de zona, sin que se genere radicado ni seguimiento alguno.

El presente proyecto propone digitalizar ese proceso mediante dos productos: una aplicación móvil con funcionamiento autónomo sin conexión para el operario de campo, y un portal web para el cliente. Los tres actores del sistema son el cliente residencial, el recaudador instalador y el administrador del sistema.

\subsection{Requerimientos Funcionales}

Los requerimientos funcionales se identificaron a partir de la observación directa del proceso operativo y de entrevistas con el personal de campo y clientes activos de CEDENAR.

\begin{longtable}{p{1.2cm} p{3cm} p{9cm} p{1.5cm}}
  \caption{Requerimientos Funcionales del Sistema} \\
  \toprule
  \textbf{ID} & \textbf{Nombre} & \textbf{Descripción} & \textbf{Prior.} \\
  \midrule
  \endfirsthead
  \caption[]{Requerimientos Funcionales del Sistema (continuación)} \\
  \toprule
  \textbf{ID} & \textbf{Nombre} & \textbf{Descripción} & \textbf{Prior.} \\
  \midrule
  \endhead
  \bottomrule
  \endfoot

  RF-01 & Registro e inicio de sesión &
    Correo electrónico y número de cédula únicos en la base de datos.
    Contraseña almacenada mediante hash seguro, nunca en texto plano.
    Correo de confirmación automático al completar el registro.
    Mensajes de error descriptivos por campo. &
    Must \\[4pt]

  RF-02 & Consulta de saldo y pago &
    Consulta del estado de cuenta por NIU en máximo tres segundos.
    Integración con pasarela que soporte PSE, tarjeta de crédito y billeteras digitales.
    Recibo en PDF con firma digital descargable.
    Actualización del estado de pago tras confirmación de la pasarela. &
    Must \\[4pt]

  RF-03 & Reporte de fallas técnicas &
    Captura de coordenadas GPS con \textit{fallback} a \texttt{codigoContrato} cuando el permiso de ubicación no está disponible.
    Generación de radicado único y envío al módulo de despacho técnico por zona geográfica.
    Notificación al cliente con número de caso asignado. &
    Must \\[4pt]

  RF-04 & Historial de consumo &
    Gráfica comparativa de los últimos seis meses frente al mismo período del año anterior.
    Alerta visible cuando el consumo supera el promedio histórico en más del 20\%.
    Tiempo máximo de carga de tres segundos. &
    Should \\[4pt]

  RF-05 & Gestión del recorrido &
    Lista de predios asignados por jornada con datos del cliente y ubicación.
    Captura de lectura por escaneo de código QR o digitación manual.
    Fotografía del medidor como evidencia adjunta.
    Validación del valor ingresado contra el histórico con indicador semáforo.
    Envío del comprobante por WhatsApp o correo electrónico.
    Funcionamiento sin conexión con sincronización automática al recuperar señal. &
    Must \\
\end{longtable}

\subsection{Requerimientos No Funcionales}

\begin{longtable}{p{1.2cm} p{3cm} p{9.5cm}}
  \caption{Requerimientos No Funcionales del Sistema} \\
  \toprule
  \textbf{ID} & \textbf{Nombre} & \textbf{Descripción} \\
  \midrule
  \endfirsthead
  \caption[]{Requerimientos No Funcionales del Sistema (continuación)} \\
  \toprule
  \textbf{ID} & \textbf{Nombre} & \textbf{Descripción} \\
  \midrule
  \endhead
  \bottomrule
  \endfoot

  RNF-01 & Disponibilidad &
    Tiempo de actividad del 99.9\% anual.
    Los módulos de reporte de emergencias disponibles de forma continua las 24 horas los 7 días de la semana. \\[4pt]

  RNF-02 & Seguridad &
    Protocolo TLS 1.3 para datos en tránsito, AES-256 para datos en reposo.
    Autenticación de doble factor para el acceso administrativo a la base de datos. \\[4pt]

  RNF-03 & Rendimiento &
    Cualquier consulta de saldo, historial o estado de radicado debe resolverse en máximo tres segundos.
    Se aplica caché en servidor y consultas indexadas a la base de datos. \\[4pt]

  RNF-04 & Usabilidad en campo &
    La aplicación opera sin conexión y sincroniza de forma autónoma al recuperar señal.
    Elementos táctiles de tamaño adecuado para uso en exteriores.
    Acceso a la cámara del dispositivo directo, sin redirección a otra aplicación. \\[4pt]

  RNF-05 & Trazabilidad y auditoría &
    Registro de auditoría de solo lectura para todos los roles del sistema, sin posibilidad de modificación ni eliminación.
    Cada evento relevante (pago, lectura, radicado) queda registrado con el identificador del usuario y la marca de tiempo.
    Los registros se conservan mínimo cinco años según la normativa del sector eléctrico colombiano. \\
\end{longtable}

\subsection{Product Backlog}

El backlog completo del proyecto está administrado como Issues en el repositorio central de GitHub (\url{https://github.com/Jenifrutica/analisisCedenar/issues}).

\begin{table}[H]
  \caption{Product Backlog -- Historias de Usuario}
  \centering
  \begin{tabular}{p{1.2cm} p{5.5cm} p{2cm} p{2.8cm} p{1.5cm}}
    \toprule
    \textbf{Issue} & \textbf{Historia de Usuario} & \textbf{MoSCoW} & \textbf{Módulo} & \textbf{Estado} \\
    \midrule
    \#6 & HU-001: Registro de cuenta & Must & Autenticación & To Do \\
    \#5 & HU-002: Consulta y pago de factura & Must & Pagos & To Do \\
    \#4 & HU-003: Reporte de falla eléctrica & Must & Reportes & To Do \\
    \#3 & HU-004: Historial de consumo & Should & Dashboard cliente & To Do \\
    \#2 & HU-005: Registro de lectura del medidor & Must & Recaudador & To Do \\
    \#1 & HU-006: Registro y seguimiento de PQRS & Should & Atención al cliente & To Do \\
    \bottomrule
  \end{tabular}
\end{table}

\subsection{Historias de Usuario Seleccionadas}

Se presentan las tres historias de usuario más representativas del MVP, con su respectivo \textit{Definition of Done} y criterios de aceptación en notación Gherkin.

\subsubsection{HU-005: Registro de Lectura del Medidor}

\noindent\textit{Como} Personal de Campo, \textit{quiero} registrar la lectura del medidor de cada predio en mi ruta, adjuntar fotografía como evidencia y enviar el comprobante en formato digital, \textit{para} eliminar el uso de papel, reducir errores de digitación y agilizar la jornada de trabajo.

\vspace{6pt}
\noindent\textit{Definition of Done}

\begin{itemize}
  \item Lista de predios asignados visible al iniciar sesión
  \item Escaneo de código QR o de barras identifica el medidor correcto dentro de la ruta
  \item Motor de Validación de Consumo invocado al momento de guardar la lectura
  \item Instancia \texttt{Anomalia} generada con estado \texttt{PENDIENTE} cuando se detecta inconsistencia
  \item Modo \textit{offline} operativo; estado de la lectura actualizado a \textit{Sincronizada} al recibir respuesta HTTP 200 del servidor
  \item Envío del comprobante por WhatsApp y correo electrónico funcional
\end{itemize}

\vspace{6pt}
\noindent\textit{Criterios de aceptación}

\begin{enumerate}
  \item \textbf{Identificación del medidor por escaneo:} dado que el operario tiene su ruta del día cargada, cuando escanea el código QR del medidor, el sistema identifica el contador correcto y solicita digitar el valor en kWh observado en el display.
  \item \textbf{Operación sin conexión a internet:} dado que el operario se encuentra en zona sin datos, cuando registra lecturas y fotografías durante su recorrido, el sistema guarda los datos con estado \textit{Pendiente de Sincronización} y actualiza ese estado a \textit{Sincronizada} cuando la API retorna HTTP~200.
  \item \textbf{Alerta por lectura anómala:} dado que el operario ha ingresado el valor del medidor, cuando el Motor de Validación de Consumo detecta una inconsistencia, el sistema genera una instancia \texttt{Anomalia} con estado \texttt{PENDIENTE} y muestra una alerta antes de confirmar el guardado.
\end{enumerate}

\subsubsection{HU-003: Reporte de Falla o Daño Eléctrico}

\noindent\textit{Como} Cliente, \textit{quiero} reportar un corte de energía o daño en la infraestructura eléctrica desde la aplicación, \textit{para} que el equipo técnico de CEDENAR atienda la falla y yo pueda hacer seguimiento mediante un número de radicado.

\vspace{6pt}
\noindent\textit{Definition of Done}

\begin{itemize}
  \item Coordenadas GPS capturadas automáticamente; \textit{fallback} a \texttt{codigoContrato} cuando el permiso no está disponible
  \item Radicado único generado y almacenado en el módulo de mantenimiento
  \item Reporte enrutado al despacho técnico según la zona geográfica detectada
  \item Cliente notificado con el número de radicado asignado
  \item Registro de auditoría generado en cada cambio de estado del radicado
\end{itemize}

\vspace{6pt}
\noindent\textit{Criterios de aceptación}

\begin{enumerate}
  \item \textbf{Reporte con GPS activo:} dado que el cliente accede al módulo ``Reportar falla'' con GPS activo en el dispositivo, cuando completa el formulario y lo envía, el sistema captura las coordenadas automáticamente, genera un radicado único, enruta el caso al despacho técnico correspondiente y notifica al cliente.
  \item \textbf{Reporte sin GPS disponible:} dado que el cliente no tiene GPS activo o ha negado el permiso de ubicación, cuando reporta la falla, el sistema usa el \texttt{codigoContrato} para determinar la zona de despacho técnico y genera el radicado notificando al cliente de la misma forma.
\end{enumerate}

\subsubsection{HU-006: Registro y Seguimiento de PQRS}

\noindent\textit{Como} Cliente, \textit{quiero} registrar una petición, queja, reclamo o sugerencia directamente desde la aplicación, \textit{para} recibir respuesta formal sin necesidad de desplazarme a una oficina.

\vspace{6pt}
\noindent\textit{Definition of Done}

\begin{itemize}
  \item Formulario de PQRS disponible para el cliente y para el Personal de Campo
  \item Clasificación automática por tipo: Petición, Queja, Reclamo o Sugerencia
  \item Número de radicado generado y notificación enviada por correo y WhatsApp al registrar
  \item Panel administrativo con filtros por estado y por tipo de PQRS
  \item Notificación automática al cliente cuando el caso queda marcado como resuelto
\end{itemize}

\vspace{6pt}
\noindent\textit{Criterios de aceptación}

\begin{enumerate}
  \item \textbf{Cliente registra PQRS desde la aplicación:} dado que el cliente ha iniciado sesión en el sistema, cuando completa el formulario y hace clic en \textit{Enviar}, el sistema genera un número de radicado único y notifica al cliente por correo y WhatsApp con el número de caso y el tiempo estimado de respuesta.
  \item \textbf{Notificación de resolución al cliente:} dado que la PQRS ha sido marcada como resuelta, cuando el sistema actualiza el estado del radicado, envía una notificación automática por correo y WhatsApp al cliente.
\end{enumerate}

% ============================================================
\newpage
\section{Arquitectura Empresarial (ArchiMate)}
% ============================================================

\noindent El modelo de arquitectura empresarial organiza el ecosistema de la solución en tres capas bajo el estándar ArchiMate. La capa de negocio define los actores Operario de Campo y Administrador de Lecturas, con los procesos de asignación de rutas de inspección, toma física del consumo en el predio, validación de consistencia y auditoría de anomalías. La capa de aplicación aloja la Aplicación Móvil de Captura y el Dashboard Web de Administración, soportados por los servicios de Autenticación, Registro de Consumo, Motor de Validación Histórica y Alertas Automáticas. La capa de tecnología comprende los dispositivos móviles corporativos con soporte GPS, la red de comunicaciones con soporte para operación sin conexión, los servidores de aplicaciones en la nube y los motores de base de datos relacional.

Cuando el Operario de Campo ejecuta el proceso de toma de lectura (capa de negocio), este se apoya en la Aplicación Móvil (capa de aplicación), que almacena y procesa los datos en el Servidor Cloud mediante servicios HTTPS seguros (capa de tecnología). La comunicación con el Servidor Central de CEDENAR se realiza a través de un enlace dedicado VPN.

\begin{figure}[H]
  \centering
  % IMAGEN: archimate.png
  % Captura del diagrama de la sección 2 de la wiki del repositorio
  \includegraphics[width=0.9\textwidth]{archimate}
  \caption{\textit{Diagrama de Arquitectura Empresarial bajo el estándar ArchiMate.}}
\end{figure}

% ============================================================
\newpage
\section{Modelado Estático (UML 2.5.1)}
% ============================================================

\clearpage
\subsection{Diagrama de Clases}

\noindent La clase abstracta \texttt{Persona} centraliza los atributos comunes de identidad. Heredan de ella mediante relaciones de generalización las clases \texttt{Supervisor}, \texttt{Cliente} y \texttt{RecaudadorInstalador}. Un \texttt{Cliente} posee una relación de agregación con uno o más \texttt{Contador}. Cada \texttt{Contador} compone de forma estricta un historial de objetos \texttt{Lectura}, que incluye el atributo \texttt{esConsistente} y la enumeración \texttt{EstadoLectura}. Cuando la condición \texttt{esConsistente = false} se evalúa como verdadera, el sistema genera una instancia de \texttt{Anomalia} vinculada al \texttt{Supervisor} para su evaluación mediante el método \texttt{evaluarAnomalia()}. El módulo comercial integra las clases \texttt{Factura} (generada a partir de una lectura válida), \texttt{Pago}, \texttt{PQRS} y \texttt{Radicado} con captura de coordenadas GPS.

\begin{figure}[H]
  \centering
  % IMAGEN: diagrama_clases.png
  % URL de descarga: https://github.com/user-attachments/assets/9dd72c89-1f8e-4d6f-bdf6-d997c4b8bad3
  \includegraphics[width=1\textwidth]{diagrama_clases.jpg}
  \caption{\textit{Diagrama de Clases -- Modelo de Dominio del sistema CEDENAR Rural.}}
\end{figure}

\clearpage
\subsection{Diagrama de Paquetes}

\noindent La estructura de código fuente adopta el patrón Domain-Driven Design / Layered Architecture bajo el espacio de nombres \texttt{com.cedenar.rural}. Cuatro paquetes con dependencias estrictamente unidireccionales: \texttt{controller} actúa como punto de entrada REST y accede a \texttt{service}, que contiene las implementaciones de lógica de negocio (incluyendo \texttt{MotorValidacionService}) y accede a \texttt{repository}, que gestiona la persistencia mediante el patrón DAO y usa \texttt{domain}. El paquete \texttt{domain} no tiene dependencias hacia ningún otro paquete y contiene exclusivamente las entidades puras del negocio.

\begin{figure}[H]
  \centering
  % IMAGEN: diagrama_paquetes.png
  % URL de descarga: https://github.com/user-attachments/assets/b8dd4082-099d-4fb4-a904-57e9e43fc593
  \includegraphics[width=0.65\textwidth]{diagrama_de_paquetes.jpeg}
  \caption{\textit{Diagrama de Paquetes -- Estructura de Código Fuente bajo patrón DDD/Layered.}}
\end{figure}

\clearpage
\subsection{Diagrama de Componentes}

\noindent La arquitectura de software adopta el enfoque de Arquitectura Orientada a Servicios (SOA) con conectores Lollipop y Socket. Los tres frontends (Aplicación Móvil del Operario, Portal Web de Supervisión y Portal Web/App del Cliente) están desacoplados y consumen las interfaces provistas por el Backend sin acceso directo a las bases de datos. El Motor de Validación de Consumo y el Gestor de Radicados se comunican internamente con el Servicio de Notificaciones mediante flujos asíncronos. Los repositorios delegan la persistencia final al Servidor Central CEDENAR mediante JDBC/SQL.

\begin{figure}[H]
  \centering
  % IMAGEN: diagrama_componentes.png
  % URL de descarga: https://github.com/user-attachments/assets/0e1f169e-0a6e-4b4a-8bd6-b031c3cebd90
  \includegraphics[width=0.88\textwidth]{diagrama_componentes.jpeg}
  \caption{\textit{Diagrama de Componentes -- Arquitectura de Software bajo patrón SOA.}}
\end{figure}

\clearpage
\subsection{Diagrama de Despliegue}

\noindent El diagrama mapea los artefactos de software sobre tres nodos de hardware: el Dispositivo Móvil del Operario transmite datos mediante red celular o satelital con protocolo HTTPS/JSON; el PC de Escritorio del Supervisor se conecta por WAN corporativa; y el Dispositivo del Cliente accede mediante HTTPS público. El Servidor Cloud centraliza el cómputo de \textit{backend}. La integración con el Servidor Central CEDENAR usa VPN dedicada y JDBC/SQL sobre la Base de Datos Comercial. La pasarela PSE/PlaceToPay se consume como nodo externo mediante REST API sobre TLS; el sistema no almacena datos bancarios localmente.

\begin{figure}[H]
  \centering
  % IMAGEN: diagrama_despliegue.png
  % URL de descarga: https://github.com/user-attachments/assets/48694deb-a298-4e39-9936-36aafa0cb2bf
  \includegraphics[width=0.88\textwidth]{diagrama_despliegue.jpeg}
  \caption{\textit{Diagrama de Despliegue -- Infraestructura Física del sistema.}}
\end{figure}

% ============================================================
\newpage
\section{Modelado Dinámico (UML 2.5.1)}
% ============================================================

\clearpage
\subsection{Diagrama de Casos de Uso}

\begin{figure}[H]
  \centering
  % IMAGEN: diagrama_casos_uso.png
  % URL de descarga: https://github.com/user-attachments/assets/7e1e6fcf-7122-4d94-9dd8-07acfbcde390
  \includegraphics[width=0.6\textwidth]{diagrama_casos_uso.png}
  \caption{\textit{Diagrama de Casos de Uso -- Interacciones entre actores y el sistema.}}
\end{figure}

\clearpage
\subsection{Diagrama de Actividades}

\noindent El diagrama modela el flujo completo de captura de lecturas en campo mediante carriles que separan las responsabilidades entre el Operario, la Aplicación Móvil, el \textit{Backend} y el Supervisor. El flujo contempla explícitamente la ausencia de conexión: las lecturas se guardan localmente con estado pendiente y se sincronizan de forma autónoma al recuperar señal. Cuando la lectura es consistente, pasa a estado válido y se sincroniza con el servidor central. Cuando no lo es, queda en estado \textit{Observada} y se genera una alerta para revisión del Supervisor.

\begin{figure}[H]
  \centering
  % IMAGEN: diagrama_actividades.png
  % URL de descarga: https://github.com/user-attachments/assets/f7cfb4fb-c9b0-46df-b0da-be4a388855bb
  \includegraphics[width=0.92\textwidth]{diagrama_actividad.png}
  \caption{\textit{Diagrama de Actividades -- Flujo completo de captura de lecturas en campo.}}
\end{figure}

\clearpage
\subsection{Diagrama de Secuencia}

\noindent La secuencia modela la interacción cronológica durante la validación técnica de una lectura. El Operario captura la lectura en la Aplicación Móvil, que envía la solicitud al \textit{Backend}. El Motor de Validación evalúa el atributo \texttt{esConsistente}: cuando es verdadero, guarda la lectura como \texttt{VALIDADA}; cuando es falso, el Gestor de Anomalías crea una instancia \texttt{Anomalia} con estado \texttt{PENDIENTE} y notifica al Portal del Supervisor para su revisión.

\begin{figure}[H]
  \centering
  % IMAGEN: diagrama_secuencia.png
  % URL de descarga: https://github.com/user-attachments/assets/ff85b2fe-8cf2-486a-bf67-e578c346c1db
  \includegraphics[width=0.95\textwidth]{diagrama_secuencia.png}
  \caption{\textit{Diagrama de Secuencia -- Validación técnica de una lectura de medidor.}}
\end{figure}

\clearpage
\subsection{Diagrama de Máquina de Estados}

\noindent El diagrama representa el ciclo de vida completo de una \texttt{Lectura} dentro del sistema. La lectura inicia en estado \texttt{Nueva} y avanza a \texttt{En captura} cuando el operario comienza el registro. Una vez confirmados el valor y la fotografía, pasa a \texttt{Capturada local}. Sin conexión, transiciona a \texttt{Pendiente de sincronización}; con conexión, avanza directamente a \texttt{En validación}. Si la condición \texttt{esConsistente = true} se cumple, la lectura pasa a \texttt{Validada} y luego a \texttt{Sincronizada con CEDENAR}, completando su ciclo. Si \texttt{esConsistente = false}, la lectura queda en estado \texttt{Observada} y se genera una instancia \texttt{Anomalia} pendiente. Desde ese estado, el Supervisor puede aprobar, rechazar u ordenar una recaptura.

\begin{figure}[H]
  \centering
  % IMAGEN: maquina_estados.png
  % URL de descarga: https://github.com/user-attachments/assets/20c73dec-7b2b-4ff1-ab85-aa8405799c60
  \includegraphics[width=0.95\textwidth]{maquina_estados.png}
  \caption{\textit{Diagrama de Máquina de Estados -- Ciclo de vida de una Lectura.}}
\end{figure}

% ============================================================
\newpage
\section{De la Idea al Prototipo}
% ============================================================

\subsection{Design Thinking}

\subsubsection{Fase 1: Empatizar}

\noindent El equipo aplicó dos métodos de investigación. El primero fue el acompañamiento en ruta: Nicolas Diaz (integrante del equipo) acompañó al recaudador Iván Diaz durante una jornada completa de lecturas en Linares, Nariño, observando el proceso paso a paso desde la espera inicial de la terminal hasta las interrupciones por señal. El segundo fue la autoetnografía: varios integrantes del equipo son clientes activos de CEDENAR y documentaron sus propias experiencias con facturas de consumo atípico y la imposibilidad de reportar cortes por un canal oficial. Los hallazgos se triangularon con quejas en redes sociales y notas de prensa local de Nariño.

\vspace{6pt}
\noindent A partir de la investigación se construyeron dos perfiles de usuario (Tabla \ref{tab:personas}).

\begin{table}[H]
  \caption{Personas -- Perfiles de Usuario}
  \label{tab:personas}
  \centering
  \begin{tabular}{p{3cm} p{5.8cm} p{5.8cm}}
    \toprule
    \textbf{Atributo} & \textbf{Iván Diaz (Operario)} & \textbf{Doña Marta (Cliente)} \\
    \midrule
    Rol & Recaudador instalador, cubre casco urbano y veredas de Linares & Clienta residencial del casco urbano de Linares \\
    Experiencia & 17 años como recaudador en CEDENAR & Clienta desde hace más de una década \\
    Herramienta principal & Terminal portátil más impresora manual & Teléfono, paga en banco \\
    Meta & Cerrar el sector completo en el menor número de jornadas & Pagar lo que es justo y reportar problemas cuando ocurren \\
    Frustración central & Un refresco accidental de pantalla le borra 15 minutos de trabajo & No tiene canal al cual llamar cuando la factura llega más alta de lo esperado \\
    \bottomrule
  \end{tabular}
\end{table}

\vspace{6pt}
\noindent Los hallazgos consolidados se presentan en la Tabla \ref{tab:hallazgos}.

\begin{table}[H]
  \caption{Hallazgos Principales -- Fase Empatizar}
  \label{tab:hallazgos}
  \centering
  \small
  \begin{tabular}{p{1.2cm} p{7.8cm} p{2.5cm} p{2.2cm}}
    \toprule
    \textbf{ID} & \textbf{Hallazgo} & \textbf{Persona} & \textbf{Trazab.} \\
    \midrule
    HG-01 & La terminal se congela hasta 15 minutos cuando cae la señal o se refresca la pantalla por accidente & Iván Diaz & RF-05, RNF-04 \\
    HG-02 & Cuando la terminal no envía, el operario llena un formato a mano o exporta un ZIP & Iván Diaz & RF-05 \\
    HG-03 & Los sectores incompletos obligan a regresar al mismo sector al día siguiente & Iván Diaz & RF-05 \\
    HG-04 & Cargar la impresora manual todo el día suma peso y retrasa el recorrido & Iván Diaz & RF-05 \\
    HG-05 & Los recibos dejados en puertas de predios cerrados se pierden sin que el cliente tenga constancia & Ambos & RF-05 \\
    HG-06 & Para aclarar una factura alta el cliente debe ir en persona a la casa del recaudador instalador & Doña Marta & RF-03, RF-04 \\
    HG-07 & Los cortes se reportan por WhatsApp personal al técnico, sin radicado ni trazabilidad & Doña Marta & RF-03 \\
    HG-08 & El cliente no recibe comprobante digital de la lectura tomada & Ambos & RF-02, RF-05 \\
    \bottomrule
  \end{tabular}
\end{table}

\subsubsection{Fase 2: Definir}

\noindent Al cruzar los hallazgos emergen tres grupos de problemas. El primero, resiliencia en campo (HG-01, HG-02, HG-03), tiene su origen en que la terminal no puede trabajar si no hay señal y cualquier interrupción borra el progreso. La solución debe funcionar aunque la conexión caiga. El segundo, equipo del operario (HG-04, HG-05), muestra que la impresora manual genera fricciones físicas y riesgos reales de pérdida de comprobantes. El tercero, canal formal para el cliente (HG-06, HG-07, HG-08), son variantes del mismo problema: el cliente no tiene adónde acudir cuando algo falla.

\vspace{6pt}
\noindent La declaración del problema (POV) resultante es la siguiente:

\begin{quote}
Iván Diaz lleva 17 años tomando lecturas en Linares y todavía depende de que haya señal en el momento exacto para no perder el trabajo de media jornada. Necesita una herramienta que guarde el progreso localmente y lo sincronice de forma autónoma cuando vuelva la conexión, sin que un refresco accidental lo obligue a empezar de nuevo.

Doña Marta no tiene canal formal para aclarar su factura ni para reportar un corte de luz. Cuando algo falla, tiene que ir físicamente a la casa del recaudador. Necesita un canal que le genere un número de radicado y le informe el estado de su caso.
\end{quote}

\vspace{6pt}
\noindent Las preguntas \textit{How Might We} (HMW) derivadas se presentan en la Tabla \ref{tab:hmw}.

\begin{table}[H]
  \caption{Preguntas How Might We -- Fase Definir}
  \label{tab:hmw}
  \centering
  \small
  \begin{tabular}{p{1.2cm} p{9.8cm} p{2.5cm}}
    \toprule
    \textbf{ID} & \textbf{Pregunta} & \textbf{Hallazgo} \\
    \midrule
    HMW-01 & ¿Cómo logramos que la aplicación siga funcionando aunque no haya señal y no pierda el progreso si la pantalla se refresca? & HG-01 \\
    HMW-02 & ¿Cómo automatizamos el reenvío de lecturas para que el operario no tenga que llenar ningún formato a mano ni exportar archivos ZIP? & HG-02 \\
    HMW-03 & ¿Cómo reemplazamos la impresora manual y aún así entregamos un comprobante válido en el momento de la visita? & HG-04, HG-05, HG-08 \\
    HMW-04 & ¿Cómo logramos cerrar el sector completo en una sola jornada aunque la conexión sea inestable? & HG-03 \\
    HMW-05 & ¿Cómo le damos al cliente un canal formal para aclarar su factura sin que tenga que ir a la casa del recaudador? & HG-06 \\
    HMW-06 & ¿Cómo formalizamos el reporte de fallas para que el cliente reciba un radicado y pueda hacerle seguimiento? & HG-07 \\
    \bottomrule
  \end{tabular}
\end{table}

\subsubsection{Fase 3: Idear}

\noindent El equipo aplicó la técnica de Brainwriting silencioso: cinco personas, seis ideas cada una, seis minutos por pregunta HMW, sin discusión durante la escritura. Las ideas se agruparon por afinidad y se priorizaron con el método MoSCoW (Tabla \ref{tab:moscow}).

\begin{table}[H]
  \caption{Priorización MoSCoW -- Fase Idear}
  \label{tab:moscow}
  \centering
  \begin{tabular}{p{2cm} p{12cm}}
    \toprule
    \textbf{Categoría} & \textbf{Funcionalidades priorizadas} \\
    \midrule
    Must & Arquitectura \textit{offline-first} con persistencia de jornada, cola de reintentos automáticos, captura con fotografía, semáforo de validación, comprobante digital (elimina impresora), formulario de aclaración de factura con radicado, reporte de fallas con radicado, autenticación segura \\[4pt]
    Should & Exportación ZIP como respaldo de contingencia, historial de consumo con gráfica comparativa, panel de anomalías, seguimiento de radicado, pago en línea, optimización de ruta por geolocalización \\[4pt]
    Could & Alertas \textit{push} de consumo elevado, notificación automática al cliente al momento de la lectura \\[4pt]
    Won't & Medidores inteligentes con lectura remota, chatbot de atención al cliente \\
    \bottomrule
  \end{tabular}
\end{table}

\noindent Esta priorización coincide con la columna MoSCoW del backlog de la Sección 1.4, lo que confirma que Design Thinking y el levantamiento de requerimientos llegaron al mismo resultado por caminos independientes.

\paragraph{Benchmarking}
Se revisaron cuatro soluciones existentes para verificar que los problemas identificados no tienen respuesta adecuada en el mercado actual, en particular para el contexto de conectividad rural de Nariño (Tabla \ref{tab:benchmarking}).

\begin{table}[H]
  \caption{Benchmarking -- Comparación con Soluciones Existentes}
  \label{tab:benchmarking}
  \centering
  \scriptsize
  \begin{tabularx}{\textwidth}{p{2.8cm} p{2cm} X X X X p{1.4cm}}
    \toprule
    \textbf{Solución} & \textbf{Tipo} & \textbf{Offline} & \textbf{Valid.} & \textbf{Recibo campo} & \textbf{Canal cliente} & \textbf{Rural} \\
    \midrule
    Itron Mobile (EE.UU.) & App operario & Parcial & Sí & No & No & No \\
    osapiens HUB (Europa) & CMMS y app & Sí, auto-sync & Sí & No & No & No \\
    App EPM (Colombia) & App cliente & N/A & N/A & Factura mensual & Sí & No \\
    Honeywell FieldSense & App operario & No documentado & No especif. & No & No & No \\
    \textbf{Propuesta CEDENAR} & \textbf{App operario} & \textbf{Sí, nativo} & \textbf{Semáforo visual} & \textbf{WhatsApp o correo} & Iteración 2 & \textbf{Sí} \\
    \bottomrule
  \end{tabularx}
\end{table}

\noindent Ninguna solución existente combina funcionamiento \textit{offline-first} nativo con generación de comprobante digital al cliente en el momento de la lectura, diseñado para conectividad rural intermitente en municipios colombianos.

\subsubsection{Fase 4: Prototipar}

\noindent Se diseñó el prototipo en Figma con alta fidelidad. La decisión de usar alta fidelidad se justifica en que las funcionalidades más críticas (comportamiento \textit{offline}, indicador de cola de sincronización, semáforo de anomalía, flujo del formulario de aclaración de factura) requieren ver el flujo completo para evaluar si funcionan correctamente.

\noindent El prototipo cubre la aplicación móvil del operario, que constituye el MVP de la primera iteración. El portal web del cliente y el panel de supervisión están especificados en los requerimientos y los modelos de las secciones anteriores, y quedan como segunda iteración.

\begin{itemize}
  \item Prototipo interactivo: \url{https://www.figma.com/proto/evDOdF75swECN4ijDjgQ6N/CEDENAR_?node-id=0-1&t=CcAxX2niYOhNCSoK-1}
  \item Archivo de diseño completo: \url{https://www.figma.com/design/evDOdF75swECN4ijDjgQ6N/CEDENAR_?node-id=0-1&t=CcAxX2niYOhNCSoK-1}
\end{itemize}

\begin{figure}[H]
  \centering
  % IMAGEN: vista_mobile.png
  % Archivo local: diagramas/vista mobile.png
  % Renombrar a vista_mobile.png y subir a Overleaf
  \includegraphics[width=0.95\textwidth]{vista_mobile.png}
  \caption{\textit{Prototipo de la Aplicación Móvil del Operario -- Vistas principales (Figma, alta fidelidad).}}
\end{figure}

\subsubsection{Fase 5: Testear}

\noindent La validación del prototipo combina dos técnicas. La primera es la trazabilidad funcional: verificar que cada historia de usuario tenga al menos una pantalla del prototipo que la satisfaga. La segunda es la evaluación heurística de Nielsen aplicada sobre las pantallas del prototipo. Los resultados de un piloto real con usuarios de CEDENAR quedan formulados como hipótesis medibles para la siguiente iteración.

\begin{table}[H]
  \caption{Matriz de Validación -- Historias de Usuario vs. Prototipo}
  \centering
  \begin{tabular}{p{1.3cm} p{6.5cm} p{2.5cm} p{2.5cm}}
    \toprule
    \textbf{HU} & \textbf{Pantalla que la satisface} & \textbf{Estado} & \textbf{Nota} \\
    \midrule
    HU-01 & Registro y login (app móvil) & Cumple & MVP \\
    HU-02 & Portal cliente: factura y pago & Pendiente & Iteración 2 \\
    HU-03 & Portal cliente: formulario y radicado & Pendiente & Iteración 2 \\
    HU-04 & Portal cliente: gráfica de consumo & Pendiente & Iteración 2 \\
    HU-05 & App móvil: QR, fotografía y semáforo & Cumple & MVP \\
    HU-06 & Portal cliente: PQRS y radicado & Pendiente & Iteración 2 \\
    \bottomrule
  \end{tabular}
\end{table}

\noindent Las hipótesis formuladas para el piloto real son las siguientes:

\begin{itemize}
  \item Si la aplicación opera con arquitectura \textit{offline-first}, los tiempos muertos por caídas de señal o refrescos accidentales caerán de aproximadamente 15 minutos por incidente a menos de 30 segundos.
  \item Si la cola de sincronización opera de forma autónoma, el operario dejará de llenar formatos a mano o exportar archivos ZIP en al menos el 90\% de las jornadas, y cerrará el sector completo en un solo día.
  \item Si el comprobante se envía de forma digital, la impresora manual deja de ser necesaria y los reclamos por comprobantes no entregados desaparecen.
  \item Si el cliente puede reportar fallas con radicado, el tiempo medio de respuesta del despacho técnico cae al menos un 25\% frente al canal informal actual.
\end{itemize}

% ============================================================
\newpage
\subsection{Lógica de Negocio en Código (Kotlin)}
% ============================================================

\noindent Se eligió Kotlin porque el prototipo es una aplicación Android y Kotlin es el lenguaje oficial de Android desde 2017, reemplazando a Java en proyectos nuevos. La sintaxis de Kotlin (\texttt{data class}, \texttt{when}, \texttt{copy}) expresa la lógica de negocio sin ruido sintáctico, lo que facilita la lectura en contexto académico y hace que el código sea directamente portable a un desarrollo real sin traducción. Las tres funciones presentadas a continuación están trazadas cada una a una pantalla específica del prototipo y a un requerimiento del backlog.

El repositorio completo con el código fuente está disponible en: \url{https://github.com/Jenifrutica/analisisCedenar}

\clearpage
\subsubsection{Función 1: Validación con Semáforo Visual}

\noindent Pantalla de origen en el prototipo: pantalla de captura de lectura con indicador verde, amarillo o rojo. Trazabilidad: RF-05, HU-005, \texttt{MotorValidacionService} (Sección 3.2).

\begin{lstlisting}[caption={LecturaService.kt -- Función validarLectura()}]
fun validarLectura(nueva: Int, historico: List<Int>): Semaforo {
    if (historico.isEmpty()) return Semaforo.VERDE
    val promedio = historico.average()
    val desviacion = (nueva - promedio) / promedio * 100
    return when {
        desviacion > 50 -> Semaforo.ROJO
        desviacion > 20 -> Semaforo.AMARILLO
        else            -> Semaforo.VERDE
    }
}
\end{lstlisting}

\noindent La función compara el valor recién capturado contra el promedio del historial de ese medidor. Si el historial está vacío (medidor recién instalado), devuelve \texttt{VERDE} para no bloquear el registro. Cuando la nueva lectura supera en más del 50\% el promedio histórico, el resultado es \texttt{ROJO}, indicando una anomalía grave que deberá ser evaluada por el Supervisor. Una desviación entre el 20\% y el 50\% produce \texttt{AMARILLO}, señal de que el operario debe revisar el dato antes de confirmar. El resultado se traduce directamente en el color del indicador visual que aparece en la pantalla de captura del prototipo.

\clearpage
\subsubsection{Función 2: Guardado Offline con Cola de Reintentos}

\noindent Pantalla de origen en el prototipo: pantalla ``Lectura guardada localmente'' con fondo verde, que aparece cuando el estado de la lectura es \texttt{PENDIENTE\_SYNC}. Trazabilidad: RF-05, RNF-04, HU-005 criterio ``Operación sin conexión a internet''.

\begin{lstlisting}[caption={LecturaService.kt -- Funciones registrarLectura() y sincronizarPendientes()}]
fun registrarLectura(lectura: Lectura, hayConexion: Boolean): EstadoLectura {
    db.guardarLocal(lectura)

    return if (hayConexion) {
        sincronizarPendientes()
        EstadoLectura.SINCRONIZADA
    } else {
        EstadoLectura.PENDIENTE_SYNC
    }
}

fun sincronizarPendientes() {
    db.obtenerNoSincronizadas().forEach { lectura ->
        val exitoso = api.enviarLectura(lectura)
        if (exitoso) {
            db.marcarSincronizada(lectura.id)
            auditLog.registrar(
                "LECTURA_SINCRONIZADA", lectura.recaudadorId, lectura.id
            )
        }
    }
}
\end{lstlisting}

\noindent \texttt{registrarLectura()} implementa el principio \textit{offline-first}: el primer paso siempre es guardar la lectura en el almacenamiento local del dispositivo, de modo que el dato nunca se pierde aunque la aplicación se cierre o el dispositivo pierda señal. Si hay conexión en ese momento, se llama de inmediato a \texttt{sincronizarPendientes()} y el estado final es \texttt{SINCRONIZADA}. Si no hay conexión, el estado queda en \texttt{PENDIENTE\_SYNC}, que es la condición que activa la pantalla ``Lectura guardada localmente'' visible en el prototipo. La función \texttt{sincronizarPendientes()} recorre todas las lecturas pendientes y las envía al servidor. Solo cambia el estado a \texttt{SINCRONIZADA} cuando recibe confirmación HTTP~200; si el servidor no responde, la lectura permanece en la cola para el siguiente intento de WorkManager.

\clearpage
\subsubsection{Función 3: Generación de Radicado y Reporte de Falla}

\noindent Pantalla de origen en el prototipo: pantalla ``PQRS Registrado con Éxito'', donde el número \texttt{\#71906-3163-1682} visible define el formato del radicado. La función \texttt{reportarFalla()} implementa el \textit{fallback} GPS documentado en HU-003. Trazabilidad: RF-03, HU-003, HU-006, clase \texttt{Radicado} (Sección 3.1).

\begin{lstlisting}[caption={RadicadoService.kt -- Funciones generarRadicado() y reportarFalla()}]
fun generarRadicado(tipo: String, zona: String, clienteId: String): Radicado {
    val numero = "#${(10000..99999).random()}-" +
                 "${(1000..9999).random()}-" +
                 "${(1000..9999).random()}"
    auditLog.registrar("RADICADO_CREADO", clienteId, numero)
    notificaciones.notificarCliente(
        clienteId, "Su caso fue registrado. Radicado: $numero"
    )
    return Radicado(numero = numero, tipo = tipo, zona = zona)
}

fun reportarFalla(
    tipo: String,
    gps: Pair<Double, Double>?,
    clienteId: String,
    codigoContrato: String
): Radicado {
    val zona = if (gps != null) {
        determinarZonaPorGps(gps.first, gps.second)
    } else {
        determinarZonaPorContrato(codigoContrato)
    }
    return generarRadicado(tipo, zona, clienteId)
}
\end{lstlisting}

\noindent \texttt{generarRadicado()} produce el número de caso en el formato \texttt{\#XXXXX-XXXX-XXXX} visible en la pantalla ``PQRS Registrado con Éxito'' del prototipo (el número \texttt{\#71906-3163-1682} que aparece en esa pantalla sigue ese mismo esquema). Además de construir el número, la función registra el evento en el log de auditoría y envía la notificación al cliente de forma inmediata. \texttt{reportarFalla()} maneja el \textit{fallback} de ubicación: si el dispositivo tiene GPS activo, determina la zona de despacho técnico a partir de las coordenadas geográficas; si el permiso de GPS no está disponible o fue denegado por el usuario, usa el \texttt{codigoContrato} del cliente para inferir la zona desde la base de datos comercial. En ambos casos la función termina llamando a \texttt{generarRadicado()}, de modo que el cliente recibe su número de caso sin importar el método de geolocalización utilizado.

% ============================================================
\newpage
\section*{Referencias}
% ============================================================
\addcontentsline{toc}{section}{Referencias}

\noindent Repositorio del proyecto. (2026). \textit{Sistema de automatización de lecturas CEDENAR}. GitHub. \url{https://github.com/Jenifrutica/analisisCedenar}

\vspace{6pt}
\noindent Repositorio del proyecto. (2026). \textit{Wiki de documentación técnica}. GitHub. \url{https://github.com/Jenifrutica/analisisCedenar/wiki}

\vspace{6pt}
\noindent Equipo de diseño. (2026). \textit{CEDENAR\_ -- Prototipo interactivo de alta fidelidad}. Figma. \url{https://www.figma.com/proto/evDOdF75swECN4ijDjgQ6N/CEDENAR_?node-id=0-1&t=CcAxX2niYOhNCSoK-1}

\vspace{6pt}
\noindent Object Management Group. (2017). \textit{Unified Modeling Language specification version 2.5.1}. \url{https://www.omg.org/spec/UML/2.5.1}

\vspace{6pt}
\noindent The Open Group. (2019). \textit{ArchiMate 3.1 specification}. \url{https://pubs.opengroup.org/architecture/archimate3-doc/}

\vspace{6pt}
\noindent Nielsen, J. (1994). \textit{Usability engineering}. Academic Press.

\end{document}